\documentclass[10pt,a4paper]{amsart}
\usepackage[margin=19mm]{geometry}
\usepackage{amsmath,amssymb,mathtools}
\usepackage[scr=rsfs,cal=euler]{mathalfa}
\usepackage{xfrac}
\usepackage[unicode,hidelinks,bookmarks=false]{hyperref}
\newcommand{\ie}{i.e.\ }
\newcommand{\cf}{cf.\ }
\newcommand{\resp}{respectively\ }
\newcommand{\Subsection}{\subsection}
\providecommand{\nobreakdash}{\nobreak\hbox{-}}
\setlength{\emergencystretch}{2em}
\multlinegap0pt
\usepackage{mathrsfs}
\usepackage{amsopn}
\usepackage{indentfirst}
\numberwithin{equation}{section}
\newtheorem{tw}{Theorem}[section]
\newtheorem{lm}[tw]{Lemma}
\newtheorem{stw}[tw]{Proposition}
\newtheorem{wn}[tw]{Corollary}
\newtheorem{uw}[tw]{Remark}
\newtheorem{df}[tw]{Definition}
\newtheorem{ex}[tw]{Example}
\newcommand{\Cbar}{{\bar{C}}}
\newcommand{\N}{\mathbb{N}}
\newcommand{\Z}{\mathbb{Z}}
\newcommand{\R}{\mathbb{R}}
\newcommand{\C}{\mathbb{C}}
\newcommand{\Ct}{\widetilde{\mathbb{C}}}
\newcommand{\X}{\mathcal{X}}
\newcommand{\XM}{\X(M)}
\newcommand{\dist}{\operatorname{dist}}
\newcommand{\const}{\operatorname{const}}
\newcommand{\grad}{\operatorname{grad}}
\newcommand{\Ric}{\operatorname{Ric}}
\newcommand{\D}{\operatorname{D}}
\newcommand{\id}{\operatorname{id}}
\newcommand{\Span}{\operatorname{span}}
\newcommand{\DD}{\mathcal{D}}
\newcommand{\DN}{\mathcal{N}}
\newcommand{\Dh}{\mathcal{H}}
\newcommand{\Dg}{\mathcal{D}_1}
\newcommand{\tr}{\operatorname{tr}}
\newcommand{\eps}{\varepsilon}
\newcommand{\lap}{\triangle}
\newcommand{\Cniesk}{{C^\infty}}
\newcommand{\nabLC}{{\widehat{\nabla}}}
\newcommand{\Jt}{\widetilde{J}}
\newcommand{\Wz}{\widetilde{\zeta}}
\newcommand{\Wn}{\widetilde{\nabla}}
\newcommand{\Wh}{\widetilde{h_1}}
\newcommand{\Whh}{\widetilde{h_2}}
\newcommand{\Ws}{\widetilde{S}}
\newcommand{\Wt}{\widetilde{\tau _1}}
\newcommand{\Wtt}{\widetilde{\tau_2}}
\newcommand{\thetazetat}{\theta _{\widetilde{\zeta}}}
\newcommand{\Hzetat}{H_{\widetilde{\zeta}}}
\newcommand{\pszm}[1]{\frac{\partial}{\partial{#1}}}
\newcommand{\ind}{\operatorname{ind}}
\newcommand{\Rad}{\operatorname{Rad}}
\begin{document}
\title[Contact variation of almost contact pseudo-metric structures]{Contact variation of almost contact pseudo-metric structures}
\author{Aleksandra Borówka, Zuzanna Szancer}
\date{}
\maketitle
\noindent\emph{Source and licence.} Contact variation of almost contact pseudo-metric structures. Version arXiv:2608.12074v1 (CC BY 4.0). This material is distributed under the Creative Commons Attribution 4.0 International licence, http://creativecommons.org/licenses/by/4.0/, as recorded in the arXiv OAI licence metadata for this exact version and shown on the abstract page https://arxiv.org/abs/2608.12074v1. Full rights notice: licenses/arxiv-2608.12074-CC-BY-4.0.txt.\par\smallskip
\noindent\textit{Editorial scope.} Counted unit: the original \texttt{lm} environment \texttt{lm::productsplit::acx} at source lines 440--476 together with its complete proof at lines 478--524; the delivered document keeps source lines 438--525 so that every referenced label and every citation resolves inside it. Native editable source is the paper's own concatenated TeX; the AMS wrapper is editorial formatting only. No publisher class, upstream script or unrelated artwork is redistributed.
\section{Local product structure (almost contact $\times$ almost complex) and other geometric properties}


\begin{lm}\label{lm::productsplit::acx}
Let $(M_1,\varphi_1,\xi_1,\eta_1)$ be an almost contact manifold and let
$(M_2,J_2)$ be an almost complex manifold. Put
$$
M:=M_1\times M_2
$$
and define an almost contact structure $(\varphi,\xi,\eta)$ on $M$ by
$$
\varphi(X_1,X_2)=\bigl(\varphi_1X_1,J_2X_2\bigr),\qquad
\xi=(\xi_1,0),\qquad
\eta(X_1,X_2)=\eta_1(X_1),
$$
where $X_1\in TM_1$ and $X_2\in TM_2$.

Let $g$ be a pseudo-Riemannian metric on $M$ such that the canonical foliations
$$
M_1\times\{q\},\qquad \{p\}\times M_2
$$
are nondegenerate and orthogonal.
If $(M,\varphi,\xi,\eta,g,\eps)$ is an almost co-K\"ahler pseudo-metric
manifold, then there exist pseudo-Riemannian metrics $g_1$ on $M_1$ and $g_2$ on $M_2$
such that
$$
g=g_1+g_2
$$
and
$(M_1,\varphi_1,\xi_1,\eta_1,g_1,\eps)$
is an almost co-K\"ahler pseudo-metric manifold and
$(M_2,J_2,g_2)$
is an almost pseudo-K\"ahler manifold.

Moreover, if $(M,\varphi,\xi,\eta,g,\eps)$ is a co-K\"ahler pseudo-metric
manifold, then $(M_1,\varphi_1,\xi_1,\eta_1,g_1,\eps)$
is a co-K\"ahler pseudo-metric manifold and
$(M_2,J_2,g_2)$
is a pseudo-K\"ahler manifold.
\end{lm}

\begin{proof}
Since $d\eta=0$, we immediately get $d\eta_1=0$. Direct computation in local frames yields
 that $g$ is
the direct product of the pseudo-metrics $g_1:=g|_{M_1\times\{q\}}$ and $g_2:=g|_{\{p\}\times M_2}$.
\par We shall show now that $(M_1,\varphi_1,\xi_1,\eta_1,g_1,\eps)$ is an almost co-K\"ahler pseudo-metric manifold. Indeed, for $X,Y\in TM_1$ we compute
\begin{align*}
g_1(\varphi_1X,\varphi_1Y)
&=g\bigl((\varphi_1X,0),(\varphi_1Y,0)\bigr)\\
&=g\bigl(\varphi(X,0),\varphi(Y,0)\bigr)\\
&=g\bigl((X,0),(Y,0)\bigr)-\eps\,\eta(X,0)\eta(Y,0)\\
&=g_1(X,Y)-\eps\,\eta_1(X)\eta_1(Y).
\end{align*}
We also have
$$
g_1(X,\xi_1)=g\bigl((X,0),(\xi_1,0)\bigr)=\eps\,\eta(X,0)=\eps\,\eta_1(X),
$$
so $(M_1,\varphi_1,\xi_1,\eta_1,g_1,\eps)$ is an almost contact pseudo-metric manifold.
Let $\Phi_1$ be  a fundamental form for $g_1$, that is for $X,Y\in TM_1$
$$
\Phi_1(X,Y):=g_1(X,\varphi_1Y)=\Phi\bigl((X,0),(Y,0)\bigr).
$$
Therefore, for all $X,Y,Z\in TM_1$,
$$
d\Phi_1(X,Y,Z)=d\Phi\bigl((X,0),(Y,0),(Z,0)\bigr)=0.
$$
Since also $d\eta_1=0$, we conclude that $(M_1,\varphi_1,\xi_1,\eta_1,g_1,\eps)$ is an almost co-K\"ahler pseudo-metric manifold.
\par Similarly, for $U,V\in TM_2$,
\begin{align*}
g_2(J_2U,J_2V)
&=g\bigl((0,J_2U),(0,J_2V)\bigr)=g\bigl(\varphi(0,U),\varphi(0,V)\bigr)\\
&=g\bigl((0,U),(0,V)\bigr)=g_2(U,V).
\end{align*}
Thus $(M_2,J_2,g_2)$ is an almost pseudo-Hermitian manifold. If $\Omega_2$ is a fundamental form for $g_2$, then for all $U,V,W\in TM_2$
$$
d\Omega_2(U,V,W)=d\Phi\bigl((0,U),(0,V),(0,W)\bigr)=0
$$
so $(M_2,J_2,g_2)$ is an almost pseudo-K\"ahler manifold.

Finally, assume that $(M,\varphi,\xi,\eta,g,\eps)$ is co-K\"ahler. Since $d\eta=0$, normality of $(\varphi,\xi,\eta)$ is equivalent to $[\varphi,\varphi]=0$.
By straightforward computations we obtain
$$
[\varphi_1,\varphi_1]=0
\qquad\text{and}\qquad
[J_2,J_2]=0.
$$
Thus $(\varphi_1,\xi_1,\eta_1)$ is normal and $J_2$ is integrable, what completes the proof of last part.
\end{proof}


\end{document}
